\label{sec:construction}

\subsection{The Block-Group Adjacency Build Process}

Building a block-group adjacency graph for a state requires four steps:

\begin{enumerate}
\item \textbf{Download block-group shapefiles}: The Census Bureau's TIGER/Line program
  distributes block-group shapefiles at the state level for each census year. Each shapefile
  contains the polygon geometry for every block group in the state. File sizes range
  from $\sim$1 MB (Wyoming) to $\sim$85 MB (California) for the 2020 geometries.

\item \textbf{Compute shared boundaries}: For each pair of block groups, determine whether
  they share a positive-length border. Two block groups are adjacent if their polygons
  share at least one edge segment of positive length. This is computed using the
  \textsc{bisect-data} crate's \texttt{SharedBoundaryDetector}, which uses a spatial
  index (R-tree) to identify candidate pairs and then computes exact boundary lengths
  using geometric intersection.

\item \textbf{Filter water adjacency}: Block groups separated only by water bodies
  (rivers, lakes, coastal waters) are not considered adjacent for redistricting purposes.
  The TIGER/Line water area shapefiles are used to identify and remove water-crossing
  adjacency edges.

\item \textbf{Validate connectivity}: The resulting adjacency graph must be fully connected
  for each state (every block group reachable from every other). Island block groups
  (physically separated from the mainland) require special handling: they are connected
  to their nearest mainland neighbor by a pseudo-edge, with a weight of 0 to indicate
  that no actual shared boundary exists.
\end{enumerate}

The \textsc{bisect} CLI command encapsulates this process:

\begin{verbatim}
bisect fetch --year 2030 --resolution block-group \
  --state <FIPS> --workers 8
\end{verbatim}

Running without \texttt{--state} processes all 50 states sequentially. The \texttt{--workers}
flag parallelizes the spatial intersection computation within each state.

\subsection{Scale Comparison: Tract vs.\ Block Group}

\begin{table}[ht]
\centering
\caption{Adjacency Graph Scale: Tract vs.\ Block-Group Resolution}
\label{tab:scale}
\begin{tabular}{lrrrr}
\toprule
 & \multicolumn{2}{c}{Tract Resolution} & \multicolumn{2}{c}{Block-Group Resolution} \\
\cmidrule(lr){2-3} \cmidrule(lr){4-5}
Metric & National & TX & National & TX (est.) \\
\midrule
Vertices (geographic units) & 242,335 & 5,411 & $\sim$1,220,000 & $\sim$27,200 \\
Edges (adjacency pairs) & $\sim$712,000 & $\sim$16,800 & $\sim$3,580,000 & $\sim$85,000 \\
Graph file size (compressed) & 18 MB & 0.4 MB & $\sim$90 MB & $\sim$2.0 MB \\
Peak RAM during construction & $\sim$2.8 GB & $\sim$0.06 GB & $\sim$14 GB & $\sim$0.3 GB \\
Construction time (8 cores) & 4.2 hrs & 0.09 hrs & $\sim$21 hrs & $\sim$0.45 hrs \\
\midrule
50-state bisect sweep time & \multicolumn{2}{c}{87 min} & \multicolumn{2}{c}{$\sim$90 min (est.)} \\
\bottomrule
\end{tabular}
\begin{tablenotes}
\small
\item Block-group estimates based on 5$\times$ scaling from tract measurements.
  Actual times will depend on 2030 block-group geography and hardware configuration.
  Bisect sweep time at block-group resolution is similar to tract because METIS runtime
  scales with adjacency complexity (edges), not raw vertex count.
\end{tablenotes}
\end{table}

The most important observation in Table~\ref{tab:scale} is that the 50-state redistricting
sweep time at block-group resolution is estimated to be approximately the same as at
tract resolution (90 minutes). This is because the bottleneck in the bisection process
is not loading the adjacency graph but running METIS partitioning, which scales with
the number of adjacency edges cut in each bisection step. For a state with $n$ block
groups and $k$ districts, each METIS call operates on a subgraph of size $n/2^d$ at
depth $d$; the total work is $O(n \log k)$ at both resolutions, with a constant factor
determined by the edge density.

\subsection{Validation Protocol}

Each block-group adjacency graph requires validation before it can be used for
redistricting. The validation protocol has three checks:

\textbf{Connectivity check}: Run BFS from an arbitrary start block group; verify
that all block groups are reachable. Failure indicates a missing pseudo-edge
for an island block group.

\textbf{Population balance check}: Verify that the sum of block-group populations
equals the state's P.L.\ 94-171 total population (within rounding). Large discrepancies
indicate group quarters allocation errors or misaligned geography vintage.

\textbf{Boundary length check}: For the 1\% of adjacent pairs with the largest shared
boundary lengths, manually verify that the computed boundary makes geographic sense
using the Census Bureau's TIGER/Line viewer. This catches projection errors (where
two block groups in different parts of the state have an apparent boundary due to
coordinate system issues) that are rare but have caused redistricting errors in
prior cycles.

The full validation protocol is documented in \texttt{bisect-data/validation/}
and can be run as:

\begin{verbatim}
bisect label-verify <label> --year 2030 --resolution block-group
\end{verbatim}
